% Imporditud: tools/import_eksperimendid.py
% Allikas: Eesti füüsikaolümpiaadi eksperimendi ülesannete
%   kogu (Rael Kalda, 2023), ülesanne Ü54, lahendus L54.
\setAuthor{EFO žürii}
\setRound{lõppvoor}
\setYear{2009}
\setNumber{P E1}
\setDifficulty{2}
\setTopic{vedelike mehaanika}
\setVahendid{pliiats, mõõtesilinder, joogitops veega. Vee tihedus on 1,0 g/cm$^3$}
\setPunktid{10}
\setAllikas{Eksperimendi kogu Ü54, lahendus lk 56}
\setLahendusseis{toores}

\prob{Pliiats}
Määrata pliiatsi tihedus.

\hint

\solu
Asetame pliitasi mõõtesilindri nii, et ta ujuks. Fikseerime ruumala muutuse $\Delta$$V_0$. Pliiatsi mass võrdub väljatõrjutud vedeliku massiga, seega m = $\rho\Delta$$V_0$. Pliiatsi ruumala leidmisel paneme tähele, et meil vett ei ole piisavalt, et pliiatsit tervenisti uputada. Küll aga me saame teda uputada kahest erinevast otsast, märkides pliiatsil koha, milleni uputame. Mõõtesilindri abil mõõdame $\Delta$$V_1$ ja $\Delta$$V_2$. Pliiatsi ruumala seega V = $\Delta$$V_1$ + $\Delta$$V_2$. Tiheduseks saame

$\rho$ = $\rho$v $\Delta$$V_0$ . $\Delta$$V_1$ + $\Delta$$V_2$
\probend
